\GalSourcePageBoundary{066}{L0065}{37}{37}

Or je dis que de là on peut tirer la valeur de $a$. Il suffit pour cela
de chercher la solution commune à cette équation et à la proposée.
Cette solution est la seule commune, car on ne peut avoir,
par exemple,
\[
F(V,b)=0,
\]
cette équation ayant un facteur commun avec l'équation semblable,
sans quoi l'une des fonctions $\varphi(a,\ldots)$ serait égale à
l'une des fonctions $\varphi(b,\ldots)$; ce qui est contre l'hypothèse.

Il suit de là que $a$ s'exprime en fonction rationnelle de $V$, et il
en est de même des autres racines.

Cette proposition\textsuperscript{(*)} est citée sans démonstration par Abel,
dans le Mémoire posthume sur les fonctions elliptiques.

\medskip
\noindent\textsc{Lemme IV.} --- \emph{Supposons que l'on ait formé l'équation en $V$,
et que l'on ait pris l'un de ses facteurs irréductibles, en sorte
que $V$ soit racine d'une équation irréductible. Soient $V,V',V'',\ldots$
les racines de cette équation irréductible. Si $a=f(V)$
est une des racines de la proposée, $f(V')$ de même sera une racine
de la proposée.}

En effet, en multipliant entre eux tous les facteurs de la forme
$V-\varphi(a,b,c,\ldots,d)$, où l'on aura opéré sur les lettres toutes
les permutations possibles, on aura une équation rationnelle en $V$,
laquelle se trouvera nécessairement divisible par l'équation en
question; donc $V'$ doit s'obtenir par l'échange des lettres dans la
fonction $V$. Soit $F(V,a)=0$ l'équation qu'on obtient en permutant
dans $V$ toutes les lettres, hors la première. On aura donc
$F(V',b)=0$, $b$ pouvant être égal à $a$, mais étant certainement
l'une des racines de l'équation proposée; par conséquent, de
même que de la proposée et de $F(V,a)=0$ est résulté $a=f(V)$,
de même il résultera de la proposée et de $F(V',b)=0$ combinées,
la suivante $b=f(V')$.

\GalHistoricalNote{\textsuperscript{(*)} Il est remarquable que de cette proposition on peut conclure que toute
équation dépend d'une équation auxiliaire telle, que toutes les racines de cette
nouvelle équation soient des fonctions rationnelles les unes des autres; car
l'équation auxiliaire en $V$ est dans ce cas.

Au surplus, cette remarque est purement curieuse. En effet, une équation qui
a cette propriété n'est pas, en général, plus facile à résoudre qu'une autre.}
